\documentclass[11pt,a4paper]{article}
\usepackage[T1]{fontenc}
\usepackage[utf8]{inputenc}
\usepackage{lmodern}
\usepackage[margin=1in]{geometry}
\usepackage{amsmath,amssymb,amsthm,mathtools}
\usepackage{microtype}
\usepackage{enumitem}
\usepackage{booktabs,array,tabularx}
\usepackage[font=small,labelfont=bf]{caption}
\usepackage{xcolor}
\usepackage{hyperref}
\definecolor{linkblue}{HTML}{185C78}
\hypersetup{colorlinks=true,linkcolor=linkblue,citecolor=linkblue,urlcolor=linkblue,
  pdftitle={Randomized Oracle Lower Bounds for Smooth Nondual Convex Optimization},
  pdfauthor={Colombo Team}}
\setlength{\parindent}{1.2em}
\setlength{\parskip}{2pt}
\setlength{\emergencystretch}{2em}
\setlist[itemize]{leftmargin=*,itemsep=3pt,topsep=4pt}
\setlist[enumerate]{leftmargin=*,itemsep=3pt,topsep=4pt}
\newtheorem{theorem}{Theorem}[section]
\newtheorem{lemma}[theorem]{Lemma}
\newtheorem{proposition}[theorem]{Proposition}
\newtheorem{corollary}[theorem]{Corollary}
\theoremstyle{definition}
\newtheorem{definition}[theorem]{Definition}
\theoremstyle{remark}
\newtheorem{remark}[theorem]{Remark}
\newcommand{\R}{\mathbb R}
\newcommand{\E}{\mathbb E}
\newcommand{\Prb}{\mathbb P}
\newcommand{\eps}{\varepsilon}
\newcommand{\norm}[1]{\left\lVert #1\right\rVert}
\newcommand{\ip}[2]{\left\langle #1,#2\right\rangle}
\newcommand{\tO}{\widetilde O}
\newcommand{\tOmega}{\widetilde\Omega}
\newcommand{\tTheta}{\widetilde\Theta}
\DeclareMathOperator{\Exp}{Exp}
\DeclareMathOperator{\sgn}{sgn}
\DeclareMathOperator{\err}{err}
\newcolumntype{Y}{>{\raggedright\arraybackslash}X}
\newcolumntype{P}[1]{>{\raggedright\arraybackslash}p{#1}}
\title{Randomized Oracle Lower Bounds for\\Smooth Nondual Convex Optimization}
\author{Colombo Team\\[3pt]
  \small\href{mailto:colombo.2026@outlook.com}{\texttt{colombo.2026@outlook.com}}}
\date{November 23, 2026}
\begin{document}
\maketitle

\begin{abstract}
We study convex minimization over an $\ell_p$ ball when gradient regularity is
measured in an $\ell_q$ norm. For $1\le p<\min\{q,2\}$ and
$1<\kappa\le2$, we prove lower bounds against arbitrary randomized algorithms
with access to exact function values and gradients. Let
$\rho=1/p-(1/q-1/2)_+$ and $a=\kappa(1+\rho)-1$.
In dimension $d\ge32(T+1)$, a single distribution of globally admissible
objectives forces error
$\Omega_{p,q,\kappa}(LR^\kappa/((T+1)^a[\log(e(T+1))]^{\kappa-1}))$
with probability at least $4/5$ after at most $T$ oracle queries.
For fixed $2\le q<\infty$, the logarithmic factor is absent.
Combined with the upper bound of Mart\'inez-Rubio, Bullins, Guzm\'an, and
Molina, this identifies the randomized oracle exponent throughout this
nondual regime. The proof transfers a signed interpolation margin into a
smooth lower bound. A simplex regularizer gives exact support locality,
while exponential offsets permit each smooth oracle reply to be simulated
with at most three hidden-sign queries in expectation, even after arbitrary
adaptive histories. Coordinate and Hadamard constructions supply the
required margins using only $O(T)$ dimensions. Every hard objective is sampled
before the interaction and has a global minimizer in the feasible ball;
the lower bound also allows queries outside that ball. The construction
extends the exponential-shift approach of Loginov et al., whose work already
establishes the matching randomized rate for $p=1,q=2$.
\end{abstract}

\noindent\textbf{Keywords.}
Oracle complexity; convex optimization; nondual geometry; randomized lower
bounds; H\"older-continuous gradients.

\medskip
\begin{samepage}
\noindent\textbf{AI Usage.}
Nearly the entire research pipeline for this paper was carried out by
Colombo, powered by GPT-5.6 Sol.
\par
\end{samepage}

\section{Introduction}\label{sec:intro}

The geometry of a feasible set and the geometry in which an objective is
smooth need not coincide. This distinction is central to the nonstandard
oracle-complexity question posed by Guzm\'an~\cite{Guz15}: how quickly can a
convex function be minimized over $RB_p^d$ when its gradient is H\"older
continuous in $\ell_q$? In the regime $p<\min\{q,2\}$, the interaction of the
two norms leads to rates that are faster than those suggested by treating the
domain only through its radius in the smoothness norm.

Recent work of Mart\'inez-Rubio, Bullins, Guzm\'an, and
Molina~\cite{MBGM26} obtains the upper error rate
\begin{equation}\label{eq:intro-rate}
  \tO_{p,q,\kappa}\!\left(\frac{LR^\kappa}{T^a}\right),
  \qquad
  a=\kappa\left(1+\frac1p-
                 \left(\frac1q-\frac12\right)_+\right)-1,
\end{equation}
for $T\le d$. Its agreement with deterministic lower bounds leaves a separate
information question: can arbitrary randomization improve this exponent?
A deterministic resisting construction does not by itself answer that
question, because the objective in a randomized lower bound must be fixed
independently of the algorithm's realized random choices.

This paper proves that the exponent in~\eqref{eq:intro-rate} remains necessary
for randomized first-order algorithms throughout
$1\le p<\min\{q,2\}$ and $1<\kappa\le2$. The proof uses a distribution of
smooth objectives chosen before any query is made. It permits exact real
oracle answers, arbitrary adaptive queries, and unrestricted computation.
The dimension requirement is linear in the query budget.

\subsection{Contribution and relation to previous work}

Our main contribution is a smooth reduction based on a signed interpolation
condition. Given public directions $g_1,\ldots,g_m$, suppose that every sign
pattern can be realized with margin $\gamma$ by a feasible point. We show that
this condition yields a distribution of smooth functions for which optimization
requires recovering most of the hidden signs. The reduction quantitatively
controls the information in the entire exact oracle reply, including its
function value and all real-valued gradient coefficients.

The reduction combines three ingredients. A squared-norm penalty on a simplex
smooths a maximum while making sufficiently submaximal affine pieces exactly
inactive. Independent exponential offsets then allow the active information
to be exposed at a constant expected cost per query. Finally, coordinate
directions for $q\ge2$ and Hadamard directions for $q<2$ provide the margin
$\gamma=Rm^{-\rho}$, where
$\rho=1/p-(1/q-1/2)_+$. For finite $q\ge2$, a penalty adapted to the coordinate
map removes the logarithmic loss in the lower bound.

These ingredients have substantial antecedents. Smoothing by a strongly
convex dual penalty is classical~\cite{Nes05}; its use in smooth oracle lower
bounds includes Guzm\'an and Nemirovski~\cite{GN15}.
Braun, Guzm\'an, and Pokutta~\cite{BGP17} develop information-theoretic
reductions based on hidden discrete information. The signed interpolation
setup and a constant floor also appear in the nonsmooth nondual lower bounds
of Mart\'inez-Rubio and Guzm\'an~\cite[Appendix H]{MG26}.
Contreras, Guzm\'an, and Mart\'inez-Rubio~\cite[Theorem 26]{CGM25} establish
randomized smooth lower bounds in general norms using both margin and
concentration assumptions. The reduction here assumes only the stated signed
margin and constructs the additional randomness through offsets.

The closest overlapping smooth result is due to Loginov et
al.~\cite[Propositions D.7--D.8]{LG26}. For $p=1,q=2$, they obtain the
randomized lower rate $\Omega_\kappa(LR^\kappa/T^{2\kappa-1})$ in dimension
linear in $T$, using exponential shifts, conditional memorylessness, and exact
smoothing locality. We build on that mechanism. Our contribution is its
transfer to signed public directions and the resulting treatment of the full
parameter regime above. In particular, the cubic rate at
$(p,q,\kappa)=(1,2,2)$ and the exponential-shift mechanism are already present
in prior work. Table~\ref{tab:literature} separates these overlapping statements
from the general result proved here. Appendix~\ref{app:probability} compares
the probability conventions.

\subsection{Why the exponent has this form}\label{sec:roadmap}

The proof separates the rate into a geometric margin and a smoothing cost.
For $m$ hidden signs, the geometry gives
$\gamma=Rm^{-\rho}$. The exponential offsets have scale proportional to
$\gamma$, while the smoothing width is proportional to $\gamma/m$.
If $\sigma$ is the strong-convexity modulus of the simplex penalty in the
relevant norm, H\"older rescaling produces the error threshold
\begin{equation}\label{eq:rate-roadmap}
  \eps_m\asymp_\kappa
  L\gamma^\kappa\left(\frac{\sigma}{m}\right)^{\kappa-1}
  =\frac{LR^\kappa\sigma^{\kappa-1}}
         {m^{\kappa\rho+\kappa-1}}.
\end{equation}
Each oracle reply requires at most three new sign queries in expectation.
Taking $m$ proportional to $T$ therefore leaves too many signs unknown for a
successful output. The power of $T$ in~\eqref{eq:rate-roadmap} is exactly
$a=\kappa(1+\rho)-1$. A universal penalty has
$\sigma\gtrsim1/\log m$, whereas the coordinate-adapted penalty has
$\sigma=1/(q-1)$ for fixed $2\le q<\infty$.

Section~\ref{sec:results} states the model and the bounds.
Section~\ref{sec:smoothing} constructs the smooth maximum.
Sections~\ref{sec:family}--\ref{sec:reduction} prove the signed-margin
reduction, including its adaptive information estimate.
Section~\ref{sec:geometry} supplies the norm geometry and completes the main
proof. Section~\ref{sec:discussion} discusses the scope of the result.

\section{Oracle model and main results}\label{sec:results}

For $1\le r\le\infty$, write
$B_r^d=\{x\in\R^d:\norm{x}_r\le1\}$ and let $r^\star$ be its conjugate
exponent, with $1/r+1/r^\star=1$ and $1/\infty=0$.
The feasible region is $Q=RB_p^d$. We consider convex continuously
differentiable functions on a neighborhood of $Q$ satisfying
\begin{equation}\label{eq:class}
  \norm{\nabla f(x)-\nabla f(y)}_{q^\star}
  \le L\norm{x-y}_q^{\kappa-1},
  \qquad x,y\in Q,
\end{equation}
where $L,R>0$ and $1<\kappa\le2$. Every hard function constructed below
satisfies~\eqref{eq:class} on all of $\R^d$.

An exact first-order query at $x$ returns $(f(x),\nabla f(x))$.
An algorithm may choose each query and its final output using its entire
transcript and an independent random seed. We assume these rules are
measurable and impose no restriction on computation between queries.
The budget $T$ is an upper bound on the number of queries on every run.
Queries and the output are feasible in the basic model. For our globally
defined hard instances, the lower bound remains valid when queries may lie
anywhere in $\R^d$.
The output need not be one of the queried points.

Let $\err_Q(f,x)=f(x)-\min_Q f$. The randomized minimax error at confidence
$2/3$ is
\[
  e_T^{\rm rand}
  =\inf_{\mathcal A}\inf\left\{\eps\ge0:
    \sup_f\Prb_{\mathcal A}
       [\err_Q(f,\widehat x)>\eps]\le\frac13\right\},
\]
where the infimum ranges over algorithms with budget $T$ and the supremum
ranges over the stated class. For deterministic algorithms, write
$e_T^{\rm det}=\inf_{\mathcal A}\sup_f\err_Q(f,\widehat x)$.
Their inverse query complexities are denoted by
$Q_{\rm rand}(\eps;d)$ and $Q_{\rm det}(\eps;d)$.

Throughout the paper,
\begin{equation}\label{eq:rho}
  1\le p<\min\{q,2\},\qquad
  \rho=\frac1p-\left(\frac1q-\frac12\right)_+,
  \qquad a=\kappa(1+\rho)-1.
\end{equation}
All logarithms are natural. Constants with subscripts $p,q,\kappa$ may depend
on these fixed parameters. In comparisons with upper bounds,
$\widetilde O$ and $\widetilde\Theta$ suppress polylogarithmic factors in the
dimension, query budget, and accuracy; they suppress no polynomial factor in
the dimension.

\begin{theorem}[Randomized lower bound]\label{thm:main}
Fix~\eqref{eq:rho}, $L,R>0$, and an integer $T\ge1$.
Suppose $d\ge32(T+1)$. There is a probability distribution $\nu_T$ on
globally convex $C^1$ functions satisfying~\eqref{eq:class} globally such that
every randomized algorithm with at most $T$ exact first-order queries on
every run satisfies
\begin{equation}\label{eq:main}
  \Prb_{f\sim\nu_T,\mathcal A}
       [\err_Q(f,\widehat x)>\tau_T]\ge\frac45,
\end{equation}
where, for a constant $c_{p,q,\kappa}>0$,
\begin{equation}\label{eq:threshold}
\tau_T=
\begin{cases}
\displaystyle
 \frac{c_{p,q,\kappa}LR^\kappa}
      {(T+1)^{\kappa(1+1/p)-1}},
 & 2\le q<\infty,\\[10pt]
\displaystyle
 \frac{c_{p,q,\kappa}LR^\kappa}
      {(T+1)^a[\log(e(T+1))]^{\kappa-1}},
 & p<q<2\ \text{or}\ q=\infty.
\end{cases}
\end{equation}
The distribution is independent of the algorithm and its random seed.
Every objective in its support is bounded below and has a global minimizer
in $Q$. The assertion remains valid when query points are unrestricted.
\end{theorem}

The theorem is distributional: the same $\nu_T$ works for all algorithms
with the prescribed budget. Averaging~\eqref{eq:main} over the sampled
objective gives, for each randomized algorithm, a fixed admissible objective
on which its failure probability is at least $4/5$. That objective may depend
on the algorithm, but not on its realized random seed.

\begin{corollary}[Matching oracle exponent]\label{cor:matching}
In the parameter range~\eqref{eq:rho}, for $T\ge1$ and
$d\ge32(T+1)$,
\begin{equation}\label{eq:error-sandwich}
  \tau_T\le e_T^{\rm rand}\le e_T^{\rm det}
  \le\tO_{p,q,\kappa}\!\left(\frac{LR^\kappa}{(T+1)^a}\right).
\end{equation}
Consequently, write $Z=LR^\kappa/\eps$ and let
$U(\eps,d)=\tO_{p,q,\kappa}(1+Z^{1/a})$ be a query budget attaining accuracy
$\eps$ in the cited deterministic upper bound.
For sufficiently large $Z$, whenever $32(U(\eps,d)+1)\le d$,
\begin{equation}\label{eq:query}
  Q_{\rm rand}(\eps;d)=\tTheta_{p,q,\kappa}(Z^{1/a}),
  \qquad
  Q_{\rm det}(\eps;d)=\tTheta_{p,q,\kappa}(Z^{1/a}).
\end{equation}
\end{corollary}

\begin{proof}
The lower inequality is Theorem~\ref{thm:main}, and a deterministic algorithm
is a special case of a randomized one. The upper rate is
\cite[Theorem 1 and the discussion on p.~3]{MBGM26}. That work gives a
polynomial-time randomized implementation and also a deterministic algorithm
using deterministic stable selectors, for which computational efficiency is
not claimed. Both are applicable to the oracle model used here.
Inverting these error bounds gives~\eqref{eq:query} when the required budgets
are within their dimensional ranges.
\end{proof}

More explicitly, inversion of the lower bound gives
\begin{equation}\label{eq:query-lower}
Q_{\rm rand}(\eps;d)\ge
\begin{cases}
 c_{p,q,\kappa}Z^{1/a},&2\le q<\infty,\\[3pt]
 \displaystyle c_{p,q,\kappa}
 \frac{Z^{1/a}}{[\log(eZ)]^{(\kappa-1)/a}},
 &p<q<2\ \text{or}\ q=\infty,
\end{cases}
\end{equation}
provided $d$ is at least a sufficiently large constant multiple of the
displayed lower-bound budget. The condition in
Corollary~\ref{cor:matching} ensures that both budgets fit.

\begin{table}[t]
\centering
\caption{Rates in the common regime $T\ge1$, $d\ge32(T+1)$, and
$1<\kappa\le2$. Entries are normalized by $LR^\kappa$.
Each lower bound holds with failure probability at least $4/5$ under one
fixed distribution. The upper bounds are from~\cite{MBGM26}.}
\label{tab:regimes}
\small
\setlength{\tabcolsep}{5pt}
\renewcommand{\arraystretch}{1.45}
\begin{tabularx}{\textwidth}{@{}P{0.21\textwidth}P{0.23\textwidth}P{0.30\textwidth}Y@{}}
\toprule
Parameter regime & Error exponent $a$ & Randomized lower bound & Upper bound\\
\midrule
$1\le p<q<2$
& $\kappa(\frac32+\frac1p-\frac1q)-1$
& $\displaystyle\frac{c_{p,q,\kappa}}{(T+1)^a[\log(e(T+1))]^{\kappa-1}}$
& $\widetilde O((T+1)^{-a})$\\
$1\le p<2\le q<\infty$
& $\kappa(1+1/p)-1$
& $c_{p,q,\kappa}(T+1)^{-a}$
& $\widetilde O((T+1)^{-a})$\\
$1\le p<2$, $q=\infty$
& $\kappa(1+1/p)-1$
& $\displaystyle\frac{c_{p,\kappa}}{(T+1)^a[\log(e(T+1))]^{\kappa-1}}$
& $\widetilde O((T+1)^{-a})$\\
\bottomrule
\end{tabularx}
\end{table}

\begin{table}[t]
\centering
\caption{Comparison with the closest bounds. All displayed bounds apply to
exact value-and-gradient access; $a$ is defined in~\eqref{eq:rho}.
``Expected'' refers to expectation over the algorithm's random seed.
The present lower bound also implies an expected-error lower bound.}
\label{tab:literature}
\small
\setlength{\tabcolsep}{4pt}
\renewcommand{\arraystretch}{1.35}
\begin{tabularx}{\textwidth}{@{}P{0.21\textwidth}P{0.18\textwidth}P{0.265\textwidth}Y@{}}
\toprule
Work & Parameter range & Error rate after $T$ queries & Guarantee and dimension\\
\midrule
Guzm\'an~\cite{Guz15}\newline (2015), Table 1
& $1\le p<q<2$; $1<\kappa\le2$
& Lower bound\newline $\displaystyle\Omega\!\left(\frac{LR^\kappa}{T^a(\log d)^{\kappa-1}}\right)$
& Deterministic local-oracle lower bound; $T\le d$.\\
Guzm\'an~\cite{Guz15}\newline (2015), Table 1
& $1\le p<2$,\newline $2\le q\le\infty$;\newline $1<\kappa\le2$
& Lower bound\newline $\displaystyle\Omega\!\left(\frac{LR^\kappa}{T^a\min\{q,\log d\}^{\kappa-1}}\right)$
& Deterministic local-oracle lower bound; $T\le d$.\\
Mart\'inez-Rubio et al.~\cite{MBGM26}\newline (2026), Theorem 1
& $1\le p<\min\{q,2\}$; $1<\kappa\le2$
& Upper bound\newline $\widetilde O_{p,q,\kappa}(LR^\kappa/T^a)$
& Feasible queries; success $1-\delta$; $T\le d$. Deterministic oracle variant also available.\\
Loginov et al.~\cite{LG26}\newline (2026), Propositions D.7--D.8
& $p=1$, $q=2$; $1<\kappa\le2$
& Lower bound\newline $\Omega_\kappa(LR^\kappa/T^{2\kappa-1})$
& Expected error; $d\ge32T+1$. Arbitrary randomized global queries.\\
This paper, Theorem~\ref{thm:main}
& $1\le p<\min\{q,2\}$; $1<\kappa\le2$
& Lower bound\newline $\tau_T$ in~\eqref{eq:threshold}; no logarithmic loss for fixed $2\le q<\infty$
& Failure $\ge4/5$; $d\ge32(T+1)$. Arbitrary randomized global queries.\\
\bottomrule
\end{tabularx}
\end{table}

Tables~\ref{tab:regimes} and~\ref{tab:literature} compare error rates at the
same query budget. The bounds identify the exponent up to logarithmic factors;
they do not establish exact logarithmic optimality. The constant in the
finite-$q$ lower bound depends on $q$: its explicit smoothing factor is
$(q-1)^{1-\kappa}$. Thus the log-free statement is for fixed finite $q$ and
does not extend uniformly to $q=\infty$.

\section{Smoothing with exact support locality}\label{sec:smoothing}

An approximate smooth maximum is not sufficient for an exact-oracle lower
bound: an inactive affine piece could still leave a small, detectable trace
in the gradient. We use a simplex penalty whose derivatives are bounded.
The resulting optimality conditions eliminate every piece more than a
specified distance below the maximum, with no error in either the value or
the gradient.

\subsection{Squared-norm penalties on the simplex}

For an integer $m\ge1$, set $n=m+1$ and define
\[
  \Delta_n=\left\{\lambda\in\R_+^{m+1}:
                     \sum_{i=0}^m\lambda_i=1\right\}.
\]
The coordinate indexed by zero will later represent a constant affine piece.

\begin{lemma}[Properties of the penalty]\label{lem:regularizer}
For $1<r\le2$, the function
$\phi_r(\lambda)=\tfrac12\norm{\lambda}_r^2$ is
$(r-1)$-strongly convex in $\ell_r$. On $\Delta_n$ it satisfies
\begin{equation}\label{eq:penalty-bounds}
  0\le\phi_r(\lambda)\le\frac12,\qquad
  0\le\partial_i\phi_r(\lambda)\le1,\qquad
  \partial_i\phi_r(\lambda)=0\quad\text{if }\lambda_i=0.
\end{equation}
In particular, the choice
\begin{equation}\label{eq:regularizer}
  \ell=\log(en),\qquad r=1+\ell^{-1},\qquad \phi=\phi_r
\end{equation}
is strongly convex in $\ell_1$ with modulus
\begin{equation}\label{eq:sigma}
  \sigma=(r-1)n^{2/r-2}\ge\frac{e^{-2}}{\log(en)}.
\end{equation}
\end{lemma}

\begin{proof}
For positive coordinates,
\[
  \partial_i\phi_r(\lambda)
     =\norm{\lambda}_r^{2-r}\lambda_i^{r-1}.
\]
On the simplex both factors are at most one, and the formula extends
continuously to the boundary. Also $\norm{\lambda}_r\le\norm{\lambda}_1=1$,
which gives the bound on $\phi_r$.
Direct differentiation yields, for any direction $h$,
\begin{align*}
 h^\top\nabla^2\phi_r(\lambda)h
 &\ge(r-1)\norm{\lambda}_r^{2-r}
                  \sum_i\lambda_i^{r-2}h_i^2\\
 &\ge(r-1)\norm{h}_r^2.
\end{align*}
For $r<2$, the last step follows by applying H\"older's inequality to
\[
  \sum_i|h_i|^r
  =\sum_i(\lambda_i^{r-2}h_i^2)^{r/2}
              \lambda_i^{r(2-r)/2}.
\]
For $r=2$ the assertion is immediate. Integration along segments gives the
strong-convexity inequality, and continuity extends it to boundary points.
The same argument with absolute values, or approximation across coordinate
hyperplanes, gives the usual squared-norm statement on the whole space.
Since $\norm{h}_r\ge n^{1/r-1}\norm{h}_1$, the $\ell_1$ modulus is
$(r-1)n^{2/r-2}$. For~\eqref{eq:regularizer},
\[
 n^{2/r-2}
 =\exp\!\left(-\frac{2\log n}{\ell+1}\right)\ge e^{-2},
\]
which proves~\eqref{eq:sigma}.
\end{proof}

The strong convexity used here is the standard squared-norm property;
see, for example,~\cite[Section 7.2]{DSSST10}. The explicit derivative bounds
in~\eqref{eq:penalty-bounds} will provide locality.

\subsection{The smooth maximum}

Let $a_i(x)=\ip{v_i}{x}+b_i$, $0\le i\le m$, with
$\norm{v_i}_{q^\star}\le1$. For $\mu>0$, use the penalty
in~\eqref{eq:regularizer} to define
\begin{equation}\label{eq:smoothmax}
 H(x)=\max_{\lambda\in\Delta_n}
       \left\{\sum_{i=0}^m\lambda_i a_i(x)-\mu\phi(\lambda)\right\}.
\end{equation}

\begin{lemma}[Smoothness and exact elimination]\label{lem:smoothing}
The function $H$ is globally convex and continuously differentiable.
Its unique maximizing weight vector $\lambda(x)$ satisfies
\begin{align}
 \nabla H(x)&=\sum_{i=0}^m\lambda_i(x)v_i,\label{eq:grad}\\
 \norm{\nabla H(x)}_{q^\star}&\le1,\qquad
 \norm{\nabla H(x)-\nabla H(y)}_{q^\star}
       \le\frac{\norm{x-y}_q}{\mu\sigma}.\label{eq:lipgrad}
\end{align}
If $M(x)=\max_i a_i(x)$, then
\begin{equation}\label{eq:approx}
  M(x)-\frac\mu2\le H(x)\le M(x).
\end{equation}
Every index with $\lambda_i(x)>0$ obeys
\begin{equation}\label{eq:active}
  a_i(x)\ge M(x)-\mu.
\end{equation}
Consequently, deleting any indices with scores strictly below $M(x)-\mu$
changes neither $H(x)$ nor $\nabla H(x)$, provided the penalty is restricted
to the corresponding face of the original simplex.
\end{lemma}

\begin{proof}
Strong convexity makes the maximizing weights unique. The maximum
differentiation rule on the compact simplex gives~\eqref{eq:grad} and
continuous differentiability. Let $A$ be the matrix with rows $v_i^\top$.
Adding the optimality inequalities at $x$ and $y$ gives
\[
 \mu\sigma\norm{\lambda(x)-\lambda(y)}_1^2
 \le\ip{A(x-y)}{\lambda(x)-\lambda(y)}
 \le\norm{x-y}_q\norm{\lambda(x)-\lambda(y)}_1.
\]
Together with $\norm{A^\top z}_{q^\star}\le\norm{z}_1$, this proves the
Lipschitz gradient bound. The gradient norm bound follows from convexity of
the dual unit ball. The inequalities in~\eqref{eq:approx} follow from
$\phi\ge0$ and evaluation at a simplex vertex, where $\phi=1/2$.

The simplex optimality conditions give a scalar $\theta$ such that
\[
 a_i(x)-\mu\partial_i\phi(\lambda(x))\le\theta,
 \qquad\text{with equality if }\lambda_i(x)>0.
\]
At an index attaining $M(x)$, the derivative bound in
Lemma~\ref{lem:regularizer} implies $\theta\ge M(x)-\mu$.
For an active index,
$a_i(x)=\theta+\mu\partial_i\phi(\lambda(x))\ge M(x)-\mu$.
All deleted weights therefore vanish at the unique maximizer. The maximizer
belongs to the retained face, which proves exact equality of both oracle
components.
\end{proof}

This is a particular instance of dual-penalty smoothing~\cite{Nes05}.
Its relevant feature for the reduction is the exact cutoff
in~\eqref{eq:active}, rather than the approximation bound alone.

\begin{lemma}[A coordinate-adapted penalty]\label{lem:coordinate-penalty}
Suppose $2\le q<\infty$, $v_0=0$, and $v_i=s_i e_i$ for
$i\in[m]$, where $s_i\in\{-1,1\}$ and $d\ge m$.
Then every conclusion of Lemma~\ref{lem:smoothing} holds with
\begin{equation}\label{eq:qpenalty}
  \phi_q(\lambda)=\frac12\norm{\lambda}_{q^\star}^2,
  \qquad \sigma_q=q^\star-1=\frac1{q-1}
\end{equation}
in place of $\phi$ and $\sigma$.
\end{lemma}

\begin{proof}
Since $1<q^\star\le2$, Lemma~\ref{lem:regularizer} gives
$\sigma_q$-strong convexity in $\ell_{q^\star}$ and the same derivative and
vertex bounds. Here
\[
  Ax=(0,s_1x_1,\ldots,s_mx_m),\qquad
  \norm{Ax}_q\le\norm{x}_q,
  \qquad \norm{A^\top z}_{q^\star}\le\norm{z}_{q^\star}.
\]
The preceding optimality argument becomes
\[
 \mu\sigma_q\norm{\lambda(x)-\lambda(y)}_{q^\star}^2
 \le\ip{A(x-y)}{\lambda(x)-\lambda(y)}
 \le\norm{x-y}_q\norm{\lambda(x)-\lambda(y)}_{q^\star}.
\]
Thus the gradient Lipschitz constant is $1/(\mu\sigma_q)$. The approximation
and elimination statements use only~\eqref{eq:penalty-bounds}, so their
proofs are unchanged.
\end{proof}

\begin{lemma}[H\"older rescaling]\label{lem:holder}
For either of the preceding smooth maxima, let $\sigma$ denote its respective
modulus and set
\begin{equation}\label{eq:beta}
  \beta=\frac{L(\mu\sigma)^{\kappa-1}}{2^{2-\kappa}}.
\end{equation}
Then $f=\beta H$ satisfies~\eqref{eq:class} globally.
\end{lemma}

\begin{proof}
The bounded-gradient and Lipschitz estimates imply
\begin{align*}
 \norm{\nabla H(x)-\nabla H(y)}_{q^\star}
 &\le\min\left\{2,\frac{\norm{x-y}_q}{\mu\sigma}\right\}\\
 &\le2^{2-\kappa}(\mu\sigma)^{1-\kappa}
                  \norm{x-y}_q^{\kappa-1}.
\end{align*}
Multiplication by~\eqref{eq:beta} gives the required H\"older constant.
\end{proof}

\section{A hard distribution from signed interpolation}\label{sec:family}

We now isolate the geometric property needed by the information argument.
For Sections~\ref{sec:family}--\ref{sec:reduction}, $Q$ may be any nonempty
compact convex subset of $\R^d$.

\begin{definition}[Signed interpolation margin]\label{def:margin}
Public vectors $g_1,\ldots,g_m\in\R^d$ have signed interpolation margin
$\gamma>0$ on $Q$ if
\begin{equation}\label{eq:margin}
 \norm{g_i}_{q^\star}\le1\quad(i\in[m]),\qquad
 \forall s\in\{-1,1\}^m\ \exists x_s\in Q:\quad
 s_i\ip{g_i}{x_s}\le-\gamma\quad(i\in[m]).
\end{equation}
\end{definition}

This condition requires a feasible witness for every sign string, but no
formula for the constrained minimizer. The directions are known to the
algorithm; only their signs and the offsets introduced next are hidden.
Section~\ref{sec:geometry} constructs the directions explicitly for $RB_p^d$.

\subsection{Sampling and the constant floor}

Independently sample all signs and offsets before the interaction:
\begin{equation}\label{eq:random}
 S_i\stackrel{\rm iid}{\sim}\operatorname{Unif}\{-1,1\},\qquad
 E_i\stackrel{\rm iid}{\sim}\Exp(\text{mean }\eta),\qquad
 \eta=\frac{\gamma}{64},\quad
 \mu=\frac{\eta}{m}=\frac{\gamma}{64m}.
\end{equation}
Our exponential convention is $\Prb(E_i>z)=e^{-z/\eta}$ for $z\ge0$.
Define
\begin{equation}\label{eq:scores}
 c=-\frac\gamma2,\qquad
 a_0(x)=c,\qquad a_i(x)=S_i\ip{g_i}{x}-E_i\quad(i\in[m]),
\end{equation}
and let $H_{S,E}$ be the smooth maximum~\eqref{eq:smoothmax} of these scores.
The objective is
\begin{equation}\label{eq:family}
  f_{S,E}=\beta H_{S,E},
\end{equation}
with $\beta$ from~\eqref{eq:beta}. We use the universal penalty unless the
coordinate penalty in Lemma~\ref{lem:coordinate-penalty} is specified.
All penalty parameters are fixed before the interaction.
The offsets are never resampled.

The constant score $a_0=c$ has two purposes. It makes the objective globally
bounded below, and it gives every sign realization a known optimal value.
The signed witness attains that value regardless of the offsets.

\begin{lemma}[Minimum and sign recovery]\label{lem:packing}
Every objective in~\eqref{eq:family} is globally admissible, bounded below,
and has a global minimizer in $Q$. More precisely,
\begin{equation}\label{eq:minimum}
  \min_{x\in Q}H_{S,E}(x)
  =\min_{x\in\R^d}H_{S,E}(x)=c-\frac\mu2.
\end{equation}
Set $\eps_m=\beta\gamma/4$. Any point $\widehat x$ with
$\err_Q(f_{S,E},\widehat x)\le\eps_m$ recovers every sign whose offset
satisfies $E_i\le\gamma/8$ through
\begin{equation}\label{eq:decoder}
  \widehat S_i=-\sgn\ip{g_i}{\widehat x}.
\end{equation}
The number of exceptional indices satisfies
\begin{equation}\label{eq:badcount}
  B=\#\{i:E_i>\gamma/8\},\qquad \E B=me^{-8}.
\end{equation}
\end{lemma}

\begin{proof}
Lemma~\ref{lem:holder} gives global regularity. Evaluating the simplex
objective at the floor vertex gives $H_{S,E}(x)\ge c-\mu/2$ for all $x$.
At the signed witness $x_S$, every nonconstant score is at most $-\gamma$,
while $a_0=c=-\gamma/2$. Since $\mu<\gamma/2$,
Lemma~\ref{lem:smoothing} eliminates all other coordinates. Hence
$H_{S,E}(x_S)=c-\mu/2$, proving~\eqref{eq:minimum} for every realization of
the offsets.

If the objective error is at most $\eps_m$, then
$H_{S,E}(\widehat x)\le c-\mu/2+\gamma/4$.
The vertex lower bound
$H_{S,E}(\widehat x)\ge a_i(\widehat x)-\mu/2$ therefore implies
\[
  S_i\ip{g_i}{\widehat x}-E_i
  \le c+\gamma/4=-\gamma/4.
\]
For $E_i\le\gamma/8$, it follows that
$S_i\ip{g_i}{\widehat x}\le-\gamma/8<0$.
This proves~\eqref{eq:decoder}; any fixed sign convention at zero is harmless
because success excludes zero at all such indices.
Finally,
$\Prb(E_i>\gamma/8)=\exp(-\gamma/(8\eta))=e^{-8}$,
which gives~\eqref{eq:badcount}.
\end{proof}

We keep the full exponential law rather than conditioning the hard
distribution on bounded offsets. Large offsets contribute only the expected
exceptional count in~\eqref{eq:badcount}. This preserves the conditional
memorylessness required by the adaptive simulation.

\section{Exact oracle simulation under adaptivity}\label{sec:simulation}

The auxiliary information problem is recovery of
$S\in\{-1,1\}^m$ through queries that reveal one coordinate $S_i$ at a time.
The offsets are auxiliary randomness independent of $S$.
We simulate a smooth oracle reply by reading only those signs that could
influence it. The analysis charges a sign at most once during the entire
execution.

\subsection{The simulation procedure}

Let $J\subseteq[m]$ be the set of signs already read. At a query point $x$,
write
\begin{equation}\label{eq:potentials}
 y_i=\ip{g_i}{x},\qquad
 u_i=|y_i|-E_i,\qquad
 M=\max\bigl(\{c\}\cup\{S_i y_i-E_i:i\in J\}\bigr).
\end{equation}
For every unread index, $u_i$ bounds its actual score from above without using
$S_i$. Repeat the following actions:
\begin{enumerate}
\item If $J=[m]$, terminate. Otherwise, if
      $\max_{i\notin J}u_i<M-\mu$, terminate and record only
      $\mathsf{STOP}$, without recording an unread maximum or its index.
\item If the stopping condition fails, choose an unread index $i$ of largest
      potential $u_i$. Record the selected index and offset, read $S_i$, and
      add $i$ to $J$.
\item Update $M\leftarrow\max\{M,S_i y_i-E_i\}$ and repeat.
\end{enumerate}
At termination, compute the smooth maximum over the face of $\Delta_n$
supported on $\{0\}\cup J$ and return its value and gradient, multiplied by
$\beta$. A deterministic rule resolves ties; the conditional distributions
below show that ties between unread potentials have probability zero.

\begin{lemma}[Exactness of the simulation]\label{lem:exact}
For every realization, the procedure returns the exact value and gradient of
the single sampled function $f_{S,E}$ at every query point.
\end{lemma}

\begin{proof}
At termination, each unread actual score is bounded by a potential strictly
below $M-\mu$. Thus $M$ is the maximum of all actual scores, and
Lemma~\ref{lem:smoothing} eliminates all unread indices exactly.
The restricted problem uses the original penalty, with zeros in omitted
coordinates; its exponent is not recomputed for the smaller face.
\end{proof}

\subsection{The exposed transcript and conditional residuals}

The simulator is a device for analyzing information cost. Its underlying
randomness includes all sampled offsets, but the proof conditions on a
smaller, exposed transcript. This transcript contains the algorithm's random
seed, its queries and replies, and the complete loop actions
\[
   \mathsf{STOP}
   \quad\text{or}\quad
   \mathsf{REVEAL}(i,E_i,S_i).
\]
There is no intermediate transcript boundary containing only a positive
decision to continue. A reveal includes the selected potential through the
selected exact offset. A stop does not disclose the numerical maximum of
the unread potentials. The optimization algorithm itself receives only the
ordinary first-order oracle replies.

\begin{lemma}[Conditional exponential residuals]\label{lem:memoryless}
At every complete action boundary, for almost every exposed history, the
unread signs are conditionally independent fair signs. The unread offsets
have conditional product law
\begin{equation}\label{eq:posterior}
   E_i=b_i+Z_i,\qquad
   Z_i\stackrel{\rm independent}{\sim}\Exp(\text{mean }\eta),
\end{equation}
where the lower endpoints $b_i\ge0$ are determined by the history.
The residuals are independent of all unread signs.
\end{lemma}

\begin{proof}
Initially $b_i=0$, and the assertion follows from~\eqref{eq:random}.
Condition on an exposed history satisfying the assertion. The next query
$x$ is fixed by this history, which includes the algorithm's seed.

For a stop at threshold $D=M-\mu$, the new information about unread variables
is precisely
\[
  u_i<D\quad\text{for all unread }i,
  \qquad\text{equivalently}\qquad E_i>|y_i|-D.
\]
These are separate lower-bound constraints on the offsets, independent of
the unread signs. Memorylessness preserves~\eqref{eq:posterior}, replacing
$b_i$ by $\max\{b_i,|y_i|-D\}$.

For a reveal of index $j$ with potential $U$, conditioning on the selected
index and its exact offset imposes on every other unread index only
\[
  u_i\le U,
  \qquad\text{equivalently}\qquad E_i\ge|y_i|-U.
\]
The continuation condition $U\ge M-\mu$ involves only the selected value
and the previous history. Thus it places no additional constraint on the
remaining variables. Their lower endpoints become
$\max\{b_i,|y_i|-U\}$, independently. The selected sign is still fair when
it is read, and disclosing it does not change the laws of the other signs
or offsets.

Formally, the conditional density of the remaining offsets has the form
\[
  C\prod_{i\notin J}e^{-E_i/\eta}\mathbf1\{E_i\ge b_i\},
\]
which factorizes and normalizes to~\eqref{eq:posterior}.
Conditioning on an exact selected real value is understood through regular
conditional densities. The oracle reply gives no additional restriction:
by Lemma~\ref{lem:exact}, it is a deterministic function of the already
revealed variables and the query. This completes the induction.
\end{proof}

\subsection{A constant expected cost per oracle reply}

\begin{lemma}[Simulation cost]\label{lem:cost}
Conditional on any exposed history before a query, the expected number of
newly read signs is at most
\begin{equation}\label{eq:cost}
  2+\frac{m\mu}{\eta}=3.
\end{equation}
Consequently, the number $N$ of distinct signs read in at most $T$ oracle
calls satisfies $\E N\le3T$.
\end{lemma}

\begin{proof}
Call a newly read sign favorable if $S_i y_i=|y_i|$.
Conditional on its selection history, a sign is favorable with probability
$1/2$ when $y_i\ne0$ and with probability one when $y_i=0$.
The number of reads up to the first favorable sign, or earlier termination,
is therefore dominated in its tail by a geometric random variable of mean
two.

If a favorable sign is read, let $U$ be its selected potential. Its actual
score equals $U$, so the updated known maximum satisfies $M\ge U$.
By Lemma~\ref{lem:memoryless}, immediately after this reveal, every unread
$j$ has
\[
 E_j=b'_j+Z_j,\qquad b'_j\ge|y_j|-U,
 \qquad Z_j\sim\Exp(\text{mean }\eta).
\]
Hence
\begin{align*}
 \Prb(u_j\ge U-\mu\mid\text{current history})
 &=\Prb(Z_j\le |y_j|-b'_j-U+\mu\mid\text{current history})\\
 &\le1-e^{-\mu/\eta}\le\frac\mu\eta.
\end{align*}
The potential values do not change during a query. Every subsequent read
must have potential at least $M-\mu\ge U-\mu$, so the number of such reads
is bounded by the number of remaining potentials in that interval. Its
conditional expectation is at most $m\mu/\eta$.
If the query terminates before a favorable sign is reached, this second
contribution is zero. Adding the two estimates proves~\eqref{eq:cost}.
Summing conditional expectations over at most $T$ calls gives $\E N\le3T$.
\end{proof}

Repeated adaptive queries can raise the lower endpoints $b_i$, but they leave
the exponential residual scale $\eta$ unchanged. This is the feature that
makes the estimate valid throughout the interaction. The construction uses
lazy exposure of one sampled objective, rather than changes to the objective
after queries are observed.

\section{From information cost to optimization error}\label{sec:reduction}

It remains to connect the number of revealed signs to the accuracy of an
arbitrary final output. The decoding lemma requires correctness only at
indices with small offsets. The following estimate accounts for the others
without requiring independence between their count and the number of queries.

\begin{lemma}[Guessing the remaining signs]\label{lem:guess}
Let $N$ be the number of signs read during the simulation, and let $B$ be
the exceptional count in~\eqref{eq:badcount}. Any estimator based on the
terminal transcript has probability at most
\begin{equation}\label{eq:guess}
 \E\,2^{-(m-N-B)_+}
 \le\frac{\E N+\E B+1}{m+1}
\end{equation}
of correctly guessing every sign with $E_i\le\gamma/8$.
\end{lemma}

\begin{proof}
For this final decoding step, condition on the terminal exposed transcript
and on all offset values. The unread signs remain independent fair signs.
Indeed, for fixed offsets and a fixed algorithm seed, the simulation is an
adaptive decision tree that reads individual signs. Each completed leaf
depends only on the signs that it has read. The output and the set of unread
small-offset indices are fixed at that leaf; neither a potential comparison
nor a stopping decision tests an unread sign. There are at least
$(m-N-B)_+$ such indices, so their conditional probability of all being
guessed correctly is at most $2^{-(m-N-B)_+}$.

This additional conditioning is used only for the decision-tree argument.
The exponential-residual law and the estimate $\E N\le3T$ were proved with
the smaller filtration of Section~\ref{sec:simulation}. We do not assert an
exponential residual law after revealing every offset.

For integers $0\le k\le m$,
\[
  2^{k-m}\le\frac{k+1}{m+1}.
\]
To see this, note that $2^{m-k}\ge m-k+1$ and therefore
$(k+1)2^{m-k}\ge(k+1)(m-k+1)\ge m+1$.
Apply the inequality with $k=\min\{m,N+B\}$ and take expectations.
No independence assumption on $N$ and $B$ is needed.
\end{proof}

\begin{theorem}[Smooth signed-margin reduction]\label{thm:transfer}
Suppose $g_1,\ldots,g_m$ satisfy~\eqref{eq:margin} on $Q$.
Use the distribution~\eqref{eq:random}--\eqref{eq:family} with the universal
penalty and modulus~\eqref{eq:sigma}. Every randomized algorithm making at
most $T$ exact first-order queries on every run satisfies
\begin{equation}\label{eq:prob}
 \Prb\bigl[\err_Q(f_{S,E},\widehat x)\le\eps_m\bigr]
 \le\frac{3T+me^{-8}+1}{m+1},
\end{equation}
where
\begin{equation}\label{eq:epsilon-margin}
 \eps_m=
 \frac{L\gamma^\kappa\sigma^{\kappa-1}}
      {2^{5\kappa-2}m^{\kappa-1}}.
\end{equation}
All functions in the distribution satisfy~\eqref{eq:class} globally and
attain their global minimum in $Q$. Queries may lie anywhere in $\R^d$.
If $g_i=e_i$ and $2\le q<\infty$, the same assertions hold using the
coordinate penalty~\eqref{eq:qpenalty} and replacing $\sigma$ by
$\sigma_q=1/(q-1)$.
\end{theorem}

\begin{proof}
By Lemma~\ref{lem:packing}, success at error $\eps_m=\beta\gamma/4$
recovers every small-offset sign. Lemmas~\ref{lem:cost} and~\ref{lem:guess}
give
\[
 \Prb(\text{success})
 \le\frac{\E N+\E B+1}{m+1}
 \le\frac{3T+me^{-8}+1}{m+1}.
\]
Substituting $\mu=\gamma/(64m)$ into
$\eps_m=\beta\gamma/4$ gives~\eqref{eq:epsilon-margin}.
Global regularity and the existence of a feasible global minimizer were
proved in Lemma~\ref{lem:packing}. The simulation used no bound on the
query location. Its analysis includes the algorithm's independent seed in
the transcript, so the probability estimate applies directly to randomized
algorithms. In the coordinate case,
Lemma~\ref{lem:coordinate-penalty} supplies the same locality and smoothness
properties with modulus $\sigma_q$; all information estimates are unchanged.
\end{proof}

Theorem~\ref{thm:transfer} is the reusable part of the argument. Geometry
enters only through the signed margin; the information estimate is uniform
over the directions and over all query strategies. It also handles an
unqueried final output without charging an additional oracle call.

\section{Norm geometry and proof of the main theorem}\label{sec:geometry}

We construct the signed interpolation system using only $m$ coordinates.
When $d>m$, every vector is padded with zeros.

\subsection{Coordinate and Hadamard designs}

\begin{lemma}[Coordinate design for $q\ge2$]\label{lem:coords}
Let $1\le p<2$, $q\ge2$, and $d\ge m$. The choices
\[
   g_i=e_i\quad(i\in[m]),\qquad \gamma=Rm^{-1/p}
\]
satisfy~\eqref{eq:margin} on $RB_p^d$.
\end{lemma}

\begin{proof}
Each $e_i$ has dual norm one. For a sign string $s$, let
$x_s=-\gamma\sum_{i=1}^m s_i e_i$.
Then $\norm{x_s}_p=\gamma m^{1/p}=R$ and
$s_i\ip{g_i}{x_s}=-\gamma$ for every $i$.
\end{proof}

\begin{lemma}[Hadamard design for $p<q<2$]\label{lem:hadamard}
Let $1\le p<q<2$, let $m$ be a power of two, and suppose $d\ge m$.
There are directions satisfying~\eqref{eq:margin} on $RB_p^d$ with
\[
   \gamma=Rm^{-\rho},\qquad
   \rho=\frac1p-\frac1q+\frac12.
\]
\end{lemma}

\begin{proof}
Let $W$ be the Sylvester Hadamard matrix of order $m$, so its entries are
$\pm1$ and $WW^\top=W^\top W=mI_m$. For completeness, these matrices are
defined by
\[
  W_1=(1),\qquad
  W_{2m}=\begin{pmatrix}W_m&W_m\\W_m&-W_m\end{pmatrix}.
\]
Let $g_i$ be row $i$ of $m^{-1/q^\star}W$. Then
$\norm{g_i}_{q^\star}=1$.
For $s\in\{-1,1\}^m$, set
\[
   x_s=-\gamma m^{1/q^\star-1}W^\top s.
\]
Matrix multiplication gives $\ip{g_i}{x_s}=-\gamma s_i$.
Moreover, $\norm{W^\top s}_2=m$, and hence
\begin{align*}
 \norm{x_s}_p
 &\le m^{1/p-1/2}\norm{x_s}_2\\
 &=\gamma m^{1/p-1/2+1/q^\star}
 =\gamma m^\rho=R.
\end{align*}
Thus every sign pattern has a feasible witness at the claimed margin.
\end{proof}

The Hadamard construction gains the factor $m^{1/q-1/2}$ in the margin
relative to the coordinate design when $q<2$. That factor is needed to obtain
the exponent in~\eqref{eq:rho}. The construction is explicit and requires
neither random rotations nor a dimension-dependent concentration estimate.

\paragraph{Why the directions carry hidden signs.}
An unsigned composition with a Hadamard matrix does not suffice. For example,
take $p=1$ and a Hadamard matrix whose first column is identically one.
For $g_i=m^{-1/q^\star}W_i$ and $x\in RB_1^m$,
\[
 \ip{g_i}{x}\ge-Rm^{-1/q^\star},
\]
with equality for every $i$ at the single public point $x=-Re_1$.
This point minimizes every coordinatewise nondecreasing smooth maximum of
the scores $\ip{g_i}{x}-E_i$, regardless of the offsets. Hidden row signs
remove this common minimizer. The simulation in
Section~\ref{sec:simulation} accounts for exactly how those signs become
visible to the oracle.

\subsection{Completing the lower bound}

\begin{proof}[Proof of Theorem~\ref{thm:main}]
Let $m$ be the smallest power of two at least $16(T+1)$. Then
\begin{equation}\label{eq:dimension}
   16(T+1)\le m<32(T+1)\le d.
\end{equation}
Use the coordinate design if $q\ge2$ and the Hadamard design if $q<2$.
Both have $\gamma=Rm^{-\rho}$.
By Theorem~\ref{thm:transfer},
\begin{equation}\label{eq:success-constant}
 \Prb[\err_Q(f_{S,E},\widehat x)\le\eps_m]
 \le\frac{3T+me^{-8}+1}{m+1}
 \le\frac3{16}+e^{-8}<\frac15.
\end{equation}
The construction of the objective is fixed independently of the algorithm.

For $p<q<2$ or $q=\infty$, use the universal penalty. Its error threshold is
\begin{align}
 \eps_m
 &=\frac{LR^\kappa\sigma^{\kappa-1}}
          {2^{5\kappa-2}m^{\kappa\rho+\kappa-1}}
   \label{eq:exactrate}\\
 &\ge\frac{e^{-2(\kappa-1)}LR^\kappa}
 {2^{5\kappa-2}m^a[\log(e(m+1))]^{\kappa-1}}.\nonumber
\end{align}
Since $m<32(T+1)$ and $m+1\le33(T+1)$, this is at least the second
threshold in~\eqref{eq:threshold}, for example with
\[
 c_{p,q,\kappa}
 =\frac{e^{-2(\kappa-1)}}
 {2^{5\kappa-2}32^a(1+\log33)^{\kappa-1}}.
\]

For $2\le q<\infty$, use the coordinate penalty. Now $\rho=1/p$ and
$\sigma_q=1/(q-1)$, so
\begin{equation}\label{eq:finiteq-rate}
 \eps_m=
 \frac{LR^\kappa}
 {2^{5\kappa-2}(q-1)^{\kappa-1}m^{\kappa(1+1/p)-1}}.
\end{equation}
This implies the first threshold in~\eqref{eq:threshold}, for example with
$c_{p,q,\kappa}=[2^{5\kappa-2}(q-1)^{\kappa-1}32^a]^{-1}$.
Taking complements in~\eqref{eq:success-constant} proves~\eqref{eq:main}.
All remaining assertions follow from Theorem~\ref{thm:transfer}.
\end{proof}

\subsection{Example: Euclidean smoothness over the \texorpdfstring{$\ell_1$}{l1} ball}

The specialization $p=1,q=2,\kappa=2$ makes the construction especially
transparent. With $m$ as in~\eqref{eq:dimension}, take
\[
  g_i=e_i,\qquad
  \gamma=\frac Rm,\qquad
  \eta=\frac{R}{64m},\qquad
  \mu=\frac{R}{64m^2},\qquad
  \phi(\lambda)=\frac12\norm{\lambda}_2^2.
\]
Here $\sigma=1$ and $\beta=L\mu$. The hard objective is
\begin{equation}\label{eq:example}
 f_{S,E}(x)=L\mu\max_{\lambda\in\Delta_{m+1}}
 \left\{-\frac{R}{2m}\lambda_0+
     \sum_{i=1}^m\lambda_i(S_i x_i-E_i)
     -\frac\mu2\norm{\lambda}_2^2\right\}.
\end{equation}
It has an $L$-Lipschitz Euclidean gradient and a feasible global minimizer
\[
    x_S=-\frac Rm(S_1,\ldots,S_m,0,\ldots,0).
\]
The error threshold is
\begin{equation}\label{eq:cubic}
    \eps_m=\frac{LR^2}{256m^3}.
\end{equation}
The success probability at this threshold is less than $1/5$ for every
algorithm with budget $T$. Together with the known upper bound, this yields
the cubic error exponent and the query exponent $1/3$.
As noted in Section~\ref{sec:intro}, the matching randomized rate in this
special case was already obtained in~\cite{LG26}.

\section{Discussion}\label{sec:discussion}

\paragraph{Scope of the matching statement.}
Theorem~\ref{thm:main} treats the full class of convex objectives with the
stated gradient regularity in the regime $1\le p<\min\{q,2\}$.
It applies to arbitrary randomized first-order algorithms, and its dimension
requirement is $d\ge32(T+1)$. Corollary~\ref{cor:matching} combines this lower
bound with an existing upper bound; the present paper does not introduce a
new optimization algorithm. Because the dimension must accommodate the
query budget, the corollary does not describe the limit $\eps\downarrow0$
at a fixed dimension.

\paragraph{Exact information and global regularity.}
The lower bound uses the full value-and-gradient oracle without a precision
restriction. Exact elimination is what prevents unread affine pieces from
being encoded in small gradient coefficients. The floor also ensures that
every sampled function has a global minimizer inside the ball. Consequently,
allowing an algorithm to query outside the feasible set does not avoid the
information bound.

\paragraph{Remaining logarithmic factors.}
For fixed finite $q\ge2$, the coordinate map allows a penalty with a
dimension-independent modulus. At $q=\infty$, its conjugate exponent is one,
and $\tfrac12\norm{\lambda}_1^2$ is constant on the simplex, so that choice
cannot provide strong convexity. The universal penalty supplies the stated
logarithmic loss. For $p<q<2$, the Hadamard map recovers the required margin,
and the same universal penalty controls smoothness. Whether the remaining
logarithmic gaps can be removed is not settled by this argument.

\appendix
\section{A probability-convention comparison}\label{app:probability}

An expected-error lower bound alone does not generally imply a lower bound
at constant failure probability. A bounded witness for the error does permit
such a conclusion. If $0\le D\le2R$ and $\E D\ge cR$, then
\[
  \E D\le\frac{cR}{2}
         +2R\Prb\!\left(D\ge\frac{cR}{2}\right)
  \quad\Longrightarrow\quad
  \Prb\!\left(D\ge\frac{cR}{2}\right)\ge\frac c4.
\]
For the overlapping $p=1,q=2$ construction,
\cite[Sections D.5.2--D.5.3, in particular equation (38)]{LG26} uses
$D=\norm{\widehat x-x^\star}_1$ and a sufficiently small smoothing error.
The resulting comparison gives a fixed positive failure probability at the
same order of objective error.

To compare that convention with success probability $2/3$, suppose an
algorithm achieves the target error on every fixed objective with that
success probability. Run it independently $k$ times, query the objective
at each output, and return the best one. Its failure probability is at most
$3^{-k}$ and its query count is at most $k(T+1)$. For any fixed positive
failure threshold, a constant $k$ makes this probability smaller than the
threshold. Both the query budget and the required linear dimension change
only by constant factors. This explains why the two probability conventions
give the same exponent in the overlapping special case.

Conversely, the present distributional lower bound directly implies an
expected-error bound. Since a feasible output has nonnegative error,
\[
  \E_{f\sim\nu_T,\mathcal A}\err_Q(f,\widehat x)
       \ge\frac45\tau_T.
\]
Averaging again yields a fixed hard objective for each randomized algorithm.

\begin{thebibliography}{9}
\small
\setlength{\itemsep}{4pt}

\bibitem{BGP17}
G\'abor Braun, Crist\'obal Guzm\'an, and Sebastian Pokutta.
Lower bounds on the oracle complexity of nonsmooth convex optimization via
information theory.
\emph{IEEE Transactions on Information Theory}, 63(7):4709--4724, 2017.
Updated author version, arXiv:1407.5144v3, 2023.
\url{https://arxiv.org/abs/1407.5144}.

\bibitem{CGM25}
Juan Pablo Contreras, Crist\'obal Guzm\'an, and David Mart\'inez-Rubio.
Non-Euclidean high-order smooth convex optimization.
Full version, arXiv:2411.08987v2, 2025.
\url{https://arxiv.org/abs/2411.08987v2}.
Extended abstract in \emph{Proceedings of the 38th Conference on Learning
Theory}, Proceedings of Machine Learning Research 291:1330, 2025.

\bibitem{DSSST10}
John C. Duchi, Shai Shalev-Shwartz, Yoram Singer, and Ambuj Tewari.
Composite objective mirror descent.
In \emph{Proceedings of the 23rd Annual Conference on Learning Theory},
pages 14--26, 2010.
\url{https://www.ambujtewari.com/research/duchi10composite.pdf}.

\bibitem{Guz15}
Crist\'obal Guzm\'an.
Open problem: The oracle complexity of smooth convex optimization in
nonstandard settings.
In \emph{Proceedings of the 28th Conference on Learning Theory},
Proceedings of Machine Learning Research 40:1761--1763, 2015.
\url{https://proceedings.mlr.press/v40/Guzman15.html}.

\bibitem{GN15}
Crist\'obal Guzm\'an and Arkadi Nemirovski.
On lower complexity bounds for large-scale smooth convex optimization.
\emph{Journal of Complexity}, 31(1):1--14, 2015.
\url{https://arxiv.org/abs/1307.5001}.

\bibitem{LG26}
Timofei Loginov, Alexander Gasnikov, Yuriy Dorn, Aleksandr Shestakov,
Nazarii Tupitsa, Osman Osmanov, and Darina Dvinskikh.
The complexity of convex optimization with mismatched geometry.
arXiv:2609.32910v1, 2026.
\url{https://arxiv.org/abs/2609.32910v1}.

\bibitem{MBGM26}
David Mart\'inez-Rubio, Brian Bullins, Crist\'obal Guzm\'an, and Mathieu Molina.
Near-optimal acceleration for smooth $\ell_p/\ell_q$ nondual convex
first-order oracle optimization.
arXiv:2609.21880v1, 2026.
\url{https://arxiv.org/abs/2609.21880v1}.

\bibitem{MG26}
David Mart\'inez-Rubio and Crist\'obal Guzm\'an.
Stable movement for nondual Lipschitz convex optimization: Efficiency and
nearly optimal oracle rates.
arXiv:2609.20701v2, 2026.
\url{https://arxiv.org/abs/2609.20701v2}.

\bibitem{Nes05}
Yurii Nesterov.
Smooth minimization of non-smooth functions.
\emph{Mathematical Programming}, 103(1):127--152, 2005.
\href{https://doi.org/10.1007/s10107-004-0552-5}{doi:10.1007/s10107-004-0552-5}.

\end{thebibliography}
\end{document}


